\section{A finite natural-rank cutoff for the remaining characteristic-three symplectic core}
\label{sec:symplectic-three-cutoff}

\paragraph{Status and scope.}
This new derivation has passed whole-proof reviews by /root and
/root/imprimitive, with independent exact arithmetic recorded separately.
The first
lemma is an all-rank elementary support estimate, including every
order-three unipotent. The final cutoff additionally uses the published
Larsen--Tiep estimate with its previously checked explicit constant.
It does not classify the nonselfdual representations at the remaining
smaller ranks. Put
\[
 S=\operatorname{Sp}_{2n}(3),\qquad n\ge4,\qquad
 \nu(T)=a-\max_\lambda\dim\ker(T-\lambda I)
 \tag{HOS3C.1}
\]
for an operator on a complex space of dimension $a$.
All representations below are ordinary complex representations.

\begin{lemma}[D-HOS-13: every element of order three]
\label{D-HOS-13}
For every nontrivial irreducible representation $\rho$ of $S$ of
degree $a$, and every element $h\in S$ of order three,
\[
                         \nu(\rho(h))\ge2a/9.
 \tag{HOS3C.2}
\]
No selfduality or Weil-module hypothesis is imposed.
\end{lemma}

\begin{proof}
We first treat a transvection. Fix a decomposition of the natural
symplectic space into paired Lagrangians $X\oplus Y$, each of dimension
$n$. The upper Siegel group
\[
 U=\{u_C:(x,y)\mapsto(x+Cy,y)\mid C=C^t\in M_n(\mathbf F_3)\}
 \tag{HOS3C.3}
\]
is abelian. It and its opposite generate $S$. Here are explicit
elementary details for this assertion. They contain the transvections
in every $e_i$ and $f_i$ direction. The usual three-transvection
product in $\langle e_i,f_i\rangle$ interchanges those two directions
up to sign. Conjugating an upper symmetric off-diagonal elementary
shear by this interchange in one coordinate gives a Levi elementary
matrix $I+cE_{ij}$, paired with its inverse transpose on $Y$.
The square of a coordinate interchange gives the Levi diagonal with
one entry $-1$. Elementary matrices and these signs generate
$\operatorname{GL}_n(3)$, so the generated group contains that Levi.
For any $y\ne0$ and any $x$, a symmetric $C$ with $Cy=-x$ exists:
change the paired coordinates so $y$ is the first coordinate vector,
prescribe the first column and row of a symmetric matrix, and change
back by congruence. Thus the shears, Levi, and coordinate interchanges
are transitive on all nonzero natural vectors. Conjugating a long-root
transvection then gives all symplectic transvections. The elementary
transvection generation proof in D-HOS-07 gives the assertion.

The opposite groups are conjugate. Their fixed spaces in $\rho$
therefore have the same dimension, and their intersection is zero,
since $\rho$ is a nontrivial irreducible $S$-module. Hence
\[
                             t:=\dim\rho^U\le a/2.
 \tag{HOS3C.4}
\]
Write $\omega=\exp(2\pi i/3)$. The dual of $U$ consists of the
characters $C\mapsto\omega^{\operatorname{tr}(BC)}$, with $B$
symmetric; the trace pairing is nondegenerate since two is invertible.
The Levi permutes these weights by congruence, so their multiplicities
are constant on every congruence orbit. For the transvection
$u_{cE_{11}}$, $c\ne0$, the distribution of a nonzero weight orbit
on the three eigenvalues equals the distribution of
$c\,v^tBv$ as $v$ runs through the nonzero vectors of $\mathbf F_3^n$.
Indeed, a uniformly chosen change-of-basis matrix has its relevant
row uniformly distributed on those vectors; quotienting by the
stabilizer of $B$ gives the same uniform orbit measure.

Every nonzero symmetric form can be diagonalized. A rank-one form
has value proportions $(1,2,0)/3$, up to interchange of the nonzero
values, when all vectors are counted. For rank at least two, two
nonzero diagonal coordinates have value proportions either
$(5,2,2)/9$ or $(1,4,4)/9$, again up to interchange. Adding the
remaining independent coordinates convolves these distributions
and cannot increase their largest entry. It follows, after removing
the zero vector, that the zero-value proportion of any nonzero form
is at most $5/9$, and every nonzero-value proportion is at most
\[
                 \frac{2\cdot3^{n-1}}{3^n-1}\le\frac79.
 \tag{HOS3C.5}
\]
For the zero proportion, removal changes a numerator at most
$5\cdot3^{n-2}$ to at most $5\cdot3^{n-2}-1$, which only improves
the bound. The rank-one case is smaller. For a nonzero value,
the rank-one bound $2\cdot3^{n-1}$ also dominates the rank-at-least-two
bound, and the last inequality follows from $3^n\ge7$.

Thus the $1$-eigenspace of a transvection has dimension at most
$t+(5/9)(a-t)\le7a/9$. Either other eigenspace has dimension at
most $(7/9)(a-t)\le7a/9$. This proves (HOS3C.2) for every
transvection, including both nonzero root parameters.

We now record a full-line Heisenberg estimate allowing the phases
which occur when a fixed line has no fixed paired vector. Let $x\in S$
have prime order $\ell$, fix a nonzero natural vector $e$, and act
on $e^\perp/\langle e\rangle$ as $g$ of rank
$r_0=\operatorname{rank}(g-I)>0$. Suppose every nontrivial power
$g^j$, $1\le j<\ell$, has that same moving rank. In the full line
radical $Q$, write its elements as $(v,z)$ with
$v\in\mathbf F_3^{2n-2}$ and $z\in\mathbf F_3$, using the
Heisenberg law of D-HOS-01. Conjugation by $x$ has the form
\[
                   (v,z)\longmapsto(gv,z+\lambda(v))
 \tag{HOS3C.6}
\]
for a linear functional $\lambda$: the center is fixed, and a
center-fixing Heisenberg automorphism with given symplectic quotient
differs from the standard lift by an inner automorphism.
This last assertion is also immediate from the law: its central
correction is additive, and the nondegenerate alternating pairing
represents every additive functional over $\mathbf F_3$.

For either nontrivial central character, let $W$ be its unique
irreducible carrier, of dimension $N=3^{n-1}$. Its Weyl operators
$L_v$ form a basis of $\operatorname{End}(W)$, as proved in D-HOS-01.
Choose an intertwiner $T_W$ for (HOS3C.6). It conjugates $L_v$ to
a unit scalar times $L_{gv}$. Choose a $Q$-invariant positive
Hermitian inner product on $W$. The operator $T_W^*T_W$ centralizes
$Q$ and is a positive scalar by Schur's lemma; dividing by its
positive square root makes $T_W$ unitary. Its $\ell$th power is
a scalar of absolute value one, so a further unit scalar rescaling
gives $T_W^\ell=I$ without losing unitarity.
For $1\le j<\ell$, the trace on this Weyl basis is a sum of
$3^{2n-2-r_0}$ unit scalars. Since this trace also equals
$|\operatorname{tr}(T_W^j)|^2$, we obtain
\[
 |\operatorname{tr}(T_W^j)|\le N3^{-r_0/2},\qquad
 \nu(T_W)\ge(1-1/\ell)(1-3^{-r_0/2})N.
 \tag{HOS3C.7}
\]
The second assertion is exact Fourier inversion for the $\ell$th
roots of unity. On a full isotypic summand the operator is
$T_W\otimes T'$; $T'$ is a finite-order multiplicity operator.
Normalized eigenspace multiplicities convolve, so the same lower
support fraction holds on that entire summand.

All nonzero linear weights of $Q$ have the same multiplicity,
because the full $\operatorname{Sp}_{2n-2}(3)$ Levi is transitive
on them. Their fixed proportion under $g$ is
$(3^{2n-2-r_0}-1)/(3^{2n-2}-1)<3^{-r_0}$. Every other weight
orbit has length $\ell$. Thus their aggregate support fraction is
at least $(1-1/\ell)(1-3^{-r_0})$, stronger than (HOS3C.7).
The remaining summand is $\rho^Q$. D-HOS-07 proved that $Q$ and
its opposite generate $S$, so its dimension is at most $a/2$.
Adding the summands gives the universal estimate
\[
 \nu(\rho(x))\ge\tfrac12(1-1/\ell)(1-3^{-r_0/2})a.
 \tag{HOS3C.8}
\]
No phase-free lift or scalar action on the full multiplicity space
has been assumed.

Apply this with $x=h$ of order three, and write
$r=\operatorname{rank}(h-I)$ on the natural space. The rank-one
case is the transvection case already handled. If $r=2$ or $3$,
then $\dim\ker(h-I)=2n-r>n$, so this fixed space is not totally
isotropic. Choose a nondegenerate plane inside it and take $e$
in that plane. The quotient operator has $r_0=r\ge2$.
If $r\ge4$, choose any nonzero fixed vector $e$. Restricting
$h-I$ to $e^\perp$ loses at most one rank and then quotienting
by $\langle e\rangle$ loses at most one more; hence
$r_0\ge r-2\ge2$. In either case every nontrivial power of
$g$ has the same rank, because $g$ has order three and
$g^2-I=(g-I)(g+I)$ with $g+I$ invertible. Equation (HOS3C.8)
now gives $(1/2)(2/3)(2/3)a=2a/9$, as required.
\end{proof}

\begin{lemma}[D-HOS-14: every inner operator at large natural rank]
\label{D-HOS-14}
Put $N_0=10^{30}$. If $n\ge N_0$, then every noncentral
$x\in\operatorname{Sp}_{2n}(3)$ and every nontrivial irreducible
ordinary representation $\rho$ satisfy
\[
                            \nu(\rho(x))\ge2a/9.
 \tag{HOS3C.9}
\]
\end{lemma}

\begin{proof}
The projective kernel is exactly the center, by simplicity of the
projective quotient and perfectness of the covering group. Take a
power of $x$ of prime order $\ell$ modulo that center. Projective
support does not increase on taking powers or changing a scalar.
If $\ell$ is odd, a central sign makes this power $h$ have actual
order $\ell$. For $\ell=3$, D-HOS-13 applies. For $\ell=2$,
one has either $h^2=I$ or $h^2=-I$.

Suppose first that $h^2=I$. Its two eigenspaces are nonzero,
nondegenerate, and even-dimensional. Change its central sign so
that the smaller eigenspace is the $+1$ eigenspace, and choose a
nondegenerate fixed plane in it. On the remaining symplectic space
the moving rank is the dimension of the larger eigenspace, at least
$n\ge4$. Equation (HOS3C.8) with $\ell=2$ gives support at least
$(1/4)(1-3^{-2})a=2a/9$.

Next suppose $\ell\ge5$ and $h$ has a natural eigenvalue $1$ or
$-1$. After the indicated central sign choice, it has actual odd
order and a nonzero fixed symplectic subspace. Choose a fixed plane.
The quotient operator is nonidentity semisimple of order $\ell$,
so its moving rank is even and at least two; all its nontrivial
powers have the same fixed space. Equation (HOS3C.8) gives
$(1/2)(4/5)(2/3)a=4a/15>2a/9$.

All remaining cases are semisimple and have neither $1$ nor $-1$
as a natural eigenvalue. This includes $h^2=-I$. Over $\mathbf F_3$
there is no other nonzero scalar. Each irreducible factor of its
natural characteristic polynomial therefore has degree at least
two. Galois conjugate eigenvalues have equal multiplicity; hence
the largest eigenspace over the algebraic closure has dimension at
most $n$ in the natural dimension $2n$. In Larsen--Tiep's support
convention,
\[
                      \operatorname{supp}(h)/(2n)\ge1/2.
 \tag{HOS3C.10}
\]

The published uniform theorem, with the explicitly smaller constant
$\sigma_0=10^{-29}$ verified in D-UNW-01 and D-OL-01, gives
\[
 \frac{|\chi(h)|}{a}\le a^{-\sigma_0\operatorname{supp}(h)/(2n)}
 \le a^{-\sigma_0/2}
 \le3^{-(n-1)\sigma_0/2}<\frac13.
 \tag{HOS3C.11}
\]
Here D-HOS-02 gives $a\ge3^{n-1}$ and
$(N_0-1)/(2\cdot10^{29})>1$. The exact primary input is
\cite[Theorem~5.5 and complete proof, accepted-author pp.~25--26]
{UNW:LT2024}, which explicitly includes the symplectic group and
uses the natural dimension in the denominator. Its available constant
$1/(6507\cdot2^{21}\cdot10^{18})>7\cdot10^{-29}$ is larger
than the value used here. The elementary trace inequality
$|\chi(h)|\ge a-2\nu(\rho(h))$ now gives
$\nu(\rho(h))>a/3>2a/9$. This exhausts the projective-prime
cases and proves (HOS3C.9).
\end{proof}

\begin{theorem}[D-HOS-15: the remaining characteristic-three source has bounded natural rank]
\label{D-HOS-15}
Suppose a finite inducing image $J$ has a normal irreducibly acting
central product of quasisimple factors with simple quotients
$\operatorname{PSp}_{2n_i}(3)$, $n_i\ge10^{30}$, and every
factor module is nonselfdual. Then its actual six-dimensional
multiplicity variation has period-image dimension at most eight.
Different such ranks, scalar extensions and arbitrary compatible
factor permutations are allowed.

In the transitive characteristic-three symplectic-core case of
the main imaginary-quadratic sixfold problem, combining this with
D-HOS-12 gives the necessary bound
\[
                                  n<10^{30}.
 \tag{HOS3C.12}
\]
This is a finite natural-rank cutoff, not a classification or
existence assertion for the remaining ranks.
\end{theorem}

\begin{proof}
In this range the universal perfect cover is the symplectic group,
by the complete exceptional-multiplier list
\cite[Chapter~5, Table~1 and preceding paragraph, p.~35]{BI:HHM}.
Automorphisms of represented central quotients lift uniquely to
this perfect cover. The full outer group at $q=3$ has order two,
represented by nonsquare-multiplier symplectic similitudes; there
is no field or graph part
\cite[\S2.2.2, Table~2.1, author p.~14]{HES:LiXia2022}.

That nontrivial diagonal class is class-inverting on characters.
Indeed every symplectic element is a product of two anti-symplectic
involutions, by the theorem stated in
\cite[Theorem~1, author p.~1]{HOS:Vinroot2004}.
An anti-symplectic involution appearing in such a product conjugates
the element to its inverse. All anti-symplectic conjugations belong
to the same outer coset, so their actions on class functions agree.
Thus the diagonal automorphism sends $\chi$ to $\overline\chi$.
A nonselfdual irreducible character is not fixed by it. Consequently
every compatible normalizer of that factor module is a scalar times
an inner operator: after correcting its inner automorphism, this
is Schur's lemma. This argument includes every represented central
quotient and does not assume that all nonreal characters are Weil.

For a factor-preserving nonscalar operator on the total tensor
space, at least one factor is nonscalar. D-HOS-14 gives its support
fraction at least $2/9$. Tensoring the other operators convolves
normalized eigenspace multiplicities and cannot increase their
maximum. If an operator nontrivially permutes factors, the weighted
tensor-cycle argument already proved in D-HOS-06 gives support
fraction at least $1/4$, stronger than $2/9$. Nonidentity scalar
inertias move the entire total space. Therefore every genuine
inertia has moving rank at least $2e/9$, where $e$ is the total
inducing degree.

The elementary degree bound gives $e\ge3^{10^{30}-1}>9$.
The actual intermediate Hurwitz factorization has
$6=2e(h-1)+D$. Genus at least two is impossible, and in genus one
$D=6$ gives at most six genuine marked points and period rank at
most six. Rank nine would therefore force genus zero, at least
twelve genuine inertia points, and $D=2e+6$. But then
\[
                    2e+6\ge12(2e/9),\qquad e\le9,
 \tag{HOS3C.13}
\]
a contradiction. The argument uses the actual multiplicity variation
and allows its coefficient field to exceed the original field.

Finally, in a transitive core all factor modules either are all
selfdual or are all nonselfdual, since automorphic pullback and
module isomorphism preserve selfduality. The first case was excluded
in the main imaginary-character-field setting by D-HOS-12, with
its original-field-to-inducing-field argument and nontrivial
bilinear multiplier. The second case with $n\ge10^{30}$ has just
been excluded. This proves (HOS3C.12), with the stated residual scope.
\end{proof}

\paragraph{Proof and source boundary.}
The new Siegel-weight and phase-tolerant Heisenberg calculations
are supplied in full. Their input from D-HOS-01,02,07 is the
previously whole-reviewed elementary carrier, degree and generation
argument. The published inputs are exactly the uniform character
theorem, ordinary automorphism and cover descriptions, and
anti-symplectic factorization. No missing maximum-eigenspace
classification, unknown character field, or unproved normalizer
compatibility is assumed. The small-rank nonselfdual case remains
open. The even-rank Hermitian-form counting lead is separate and
is not used in this cutoff.

